% Literature Review: LLM-as-a-Judge for Hallucination Detection
% Subsection covering the paradigm, biases, and empirical validation

\subsection{LLM-as-a-Judge: Automated Evaluation of LLM Outputs}
\label{sec:lit_llm_as_judge}

The \textbf{LLM-as-a-Judge} paradigm employs large language models to evaluate the outputs of other LLMs, providing scalable automated assessment where human evaluation is prohibitively expensive or time-constrained~\citep{zheng2024judging}. This approach has emerged as a primary evaluation methodology for open-ended generation tasks where traditional metrics (BLEU, ROUGE) fail to capture semantic quality.

\subsubsection{Foundations and Empirical Validation}

The paradigm was systematically established by \citet{zheng2024judging}, who introduced MT-Bench---a multi-turn benchmark of 80 challenging questions---and Chatbot Arena, a crowdsourced platform for pairwise model comparison. Their key finding was that strong LLM judges (GPT-4) achieve over 80\% agreement with human preferences on pairwise comparisons, approaching the level of inter-annotator agreement among humans themselves. This high concordance established LLM-as-a-Judge as a viable proxy for human evaluation at scale.

Subsequent studies have confirmed and extended these findings. Recent comprehensive surveys indicate that LLM-as-a-Judge-derived accuracy metrics align more closely with human evaluations than traditional context-insensitive metrics, achieving up to 80\% agreement with human evaluators across diverse tasks~\citep{li2024llmsjudges}. This makes LLM judges strong candidates for evaluating context-aware outputs---such as causal loop diagram edges---especially when expert human evaluation is scarce.

\subsubsection{Systematic Biases in LLM Judges}

However, using one LLM to evaluate another introduces structural biases inherent to the transformer architecture and training process. Research has identified several systematic biases that affect judge reliability:

\paragraph{Position Bias.} LLM judges exhibit systematic preferences based on the \textit{order} in which candidate responses are presented. In pairwise comparison settings, judges may favor responses presented first (primacy bias) or last (recency bias), independent of actual quality~\citep{zheng2024judging}. This positional preference can flip evaluation outcomes when response order is reversed, introducing noise unrelated to content quality.

\paragraph{Verbosity Bias.} LLM judges tend to prefer longer, more detailed responses even when concise answers are equally or more accurate~\citep{zheng2024judging, li2024llmsjudges}. This bias is particularly problematic for scientific evaluation, where precision and accuracy may be better served by focused claims rather than elaborate explanations. In CLD evaluation, this could manifest as preference for edges with verbose justifications over those with succinct but accurate mechanistic explanations.

\paragraph{Self-Enhancement Bias.} LLM judges exhibit a tendency to rate outputs from their own model family more favorably than outputs from other models~\citep{li2024llmsjudges}. This self-preference introduces systematic bias when the same model (or model family) is used for both generation and evaluation, creating a form of circular validation.

\paragraph{Susceptibility to Adversarial Manipulation.} LLM-as-a-Judge has been shown to be susceptible to adversarial attacks that exploit these biases~\citep{shi2024optimization, li2024llmsjudges}. Adversarial prompts can be crafted to inflate judge scores without improving actual output quality, raising concerns about robustness in high-stakes evaluation settings.

\subsubsection{Judge Types: Pointwise vs. Pairwise Evaluation}

LLM judges can operate in two primary modes, each with distinct trade-offs:

\paragraph{Pointwise Evaluation.} The judge assesses a single response against defined criteria, producing an absolute score (e.g., 0--1 scale) or categorical verdict (correct/incorrect/partially correct). This mode is computationally efficient and enables direct comparison across different evaluation sessions, but requires well-calibrated scoring criteria.

\paragraph{Pairwise Evaluation.} The judge compares two candidate responses and selects the better one, often with explanation. This relative comparison is more robust to calibration issues---the judge need only rank, not score---but requires $O(n^2)$ comparisons for $n$ candidates and cannot produce absolute quality estimates.

For CLD hallucination detection, pointwise evaluation is more practical: each edge requires an independent assessment of validity rather than comparison against alternatives. The challenge lies in establishing reliable scoring criteria and thresholds that distinguish hallucinations from valid relationships.

\subsubsection{LLM-as-a-Corrector: Beyond Detection to Remediation}

Extending the judge paradigm, \textbf{LLM-as-a-Corrector} uses language models to not only identify errors but actively remediate them~\citep{shinn2024reflexion, dhuliawala2024chain}. In this setup, a judge first identifies problematic outputs, then a corrector LLM receives the original output along with the judge's critique and generates an improved version.

This judge-then-correct pipeline has shown promise in iterative refinement settings. Chain-of-Verification (CoVe)~\citep{dhuliawala2024chain} applies this pattern by: (1) generating a baseline response, (2) planning verification questions, (3) answering questions independently to avoid bias, and (4) generating a corrected response incorporating verification findings. Similar approaches include Reflexion~\citep{shinn2024reflexion}, which maintains a memory of past errors to guide future corrections.

However, correction introduces its own failure modes. Correctors may introduce new errors while fixing detected ones, particularly when the underlying knowledge is genuinely uncertain. In CLD contexts, a corrector might ``fix'' a valid but uncommon causal relationship by replacing it with a more stereotypical one, trading one form of error for another.

\subsubsection{Implications for CLD Hallucination Detection}

The LLM-as-a-Judge paradigm offers a scalable approach to CLD validation but requires careful consideration of its limitations:

\begin{enumerate}
    \item \textbf{Bias mitigation:} Position and verbosity biases can be partially addressed through prompt engineering (randomizing presentation order, penalizing excessive length) or ensemble methods (aggregating multiple judge calls with varied conditions).
    
    \item \textbf{Complementary validation:} Given systematic biases and the leniency pattern observed in our validation, LLM judges should be combined with context-insensitive metrics (logit-based uncertainty, embedding similarity) rather than used in isolation.
    
    \item \textbf{Conservative thresholding:} The asymmetric error profile (leniency bias) suggests that edges flagged as low-confidence by LLM judges warrant serious scrutiny, while edges receiving high scores may still benefit from additional validation.
    
    \item \textbf{Domain-specific calibration:} The observed variability in judge agreement across domains ($\kappa = 0.273$ to $0.792$) indicates that judge thresholds and interpretation may require domain-specific calibration.
\end{enumerate}

These considerations motivate the hybrid detection approach explored in this thesis, combining LLM-based judges with context-insensitive metrics to achieve more robust hallucination detection than either method alone.






